

%% ===================================================================
\section{시그모이드 함수와 바운딩}
\label{sec:lit-sigmoid}
%% ===================================================================

%% -------------------------------------------------------------------
\subsection{시그모이드 함수의 수학적 성질}
\label{subsec:sigmoid-properties}
%% -------------------------------------------------------------------

$\DRTO$ 모델에서는 관측성과 불확실성의 복합 효과를 물리적으로 타당한 범위 내로
제한하는 바운딩(bounding) 기법이 필수적이다. 복구 시간은 물리적으로
음수가 될 수 없으며, 조직이 허용하는 최대 복구 시간($R_{\max}$)을 초과할 수도 없기 때문이다.
기준 RTO($R_0$)와 최대 복구 시간 MTPD($R_{\max}$)의 상세 정의는 2.7절에서 제시한다.

\chg{본 연구의 바운딩 함수는 \citet{verhulst1838notice}가 인구 성장 모형으로 처음
도입한 표준 로지스틱 함수 $\sigma(z) = 1/(1+e^{-z})$에 이론적 뿌리를 둔다. 로지스틱
함수는 매끄러운 S자 형태로 출력을 유한 구간 안에 제한하는 성질 덕분에, 생물학적 성장
모형에서 출발하여 통계학과 여러 응용 분야의 표준적 변환으로 정착해 왔다~\citep{cramer2002origins}.}
본 연구에서 채택한 수정 시그모이드(modified sigmoid) 함수는 다음과 같이 정의된다.
\begin{equation}
  S(z) = \frac{2}{1 + e^{-z}} - 1, \quad S: \mathbb{R} \to (-1,\; 1)
  \label{eq:sigmoid-def}
\end{equation}

이 함수는 표준 로지스틱 함수를 $(-1, +1)$ 구간으로 사상한 것이다.
\chg{표준 로지스틱 함수의 치역이 $(0,1)$인 것과 달리 이를 $(-1,+1)$로 사상하여 채택한
이유는, $\DRTO$ 모델에서 관측성은 복구 시간을 기준선 \emph{아래}로, 불확실성은 기준선
\emph{위}로 작용하므로 $S(0)=0$을 중심으로 음과 양의 양방향으로 대칭인 출력이 필요하기
때문이다. 즉 로지스틱 함수의 ``매끄러운 유한 바운딩''이라는 성질은 그대로 가져오되,
두 채널의 상하 대칭 작용을 표현하기 위해 치역만 원점 대칭으로 \chg{평행이동하고 확대한} 것이다.}
%% [G19/일원화] 구 성질 열거(S(0)=0·극한·도함수·홀함수·C∞)는 2.9절 subsec:basic-properties(2계 도함수·변곡점 포함 엄밀 정리)와 중복 → 한 문장으로 축약·전방참조.
\chg{이 함수는 $S(0) = 0$의 중립성, $(-1, +1)$로의 매끄러운 유한 바운딩, 전 구간 단조증가,
$S(-z) = -S(z)$의 원점 대칭(홀함수), $C^{\infty}$ 미분 가능성 등의 핵심 성질을 가지며, 이들
성질의 엄밀한 정리와 2계 도함수와 변곡점 분석은 모델의 수학적 성질을 다루는 2.9절에서 종합한다.}


%% -------------------------------------------------------------------
\subsection{\chg{선형과 클리핑, 시그모이드 바운딩 기법의 비교}}
\label{subsec:bounding-comparison}
%% -------------------------------------------------------------------

\chg{$\DRTO$ 모델의 바운딩 함수는 다음 다섯 가지 설계 요구조건을 충족해야 한다. 첫째,
출력이 언제나 물리적으로 타당한 범위 $0 < \DRTO < R_{\max}$ 안에 머물러야 한다(경계 보장).
둘째, 입력이 작은 평균 영역에서는 근사적으로 선형으로 반응하여 해석이 직관적이어야 한다
(평균 영역 선형성). 셋째, 입력이 극단적으로 커지는 영역에서는 출력이 포화하여 경계를
넘지 않아야 한다(극단 영역 포화). 넷째, 관측성과 불확실성 두 채널의 기여가 각각의 물리적
범위 안에서 분리되어 작용해야 한다(채널 분리). 다섯째, 입력 변화에 대해 출력이 단조적으로
반응하여 입력의 순서가 보존되어야 한다(단조성). 이 다섯 요구조건을 기준으로 세 가지 후보
바운딩 기법을 비교한다.}

\textbf{선형 모델}은 $\DRTO = R_0 + \alpha \cdot x$ 형태로서, 구현이 가장 단순하고
해석이 용이하다. 그러나 입력값이 극단적인 경우 $\DRTO$가 음수가 되거나
$R_{\max}$를 초과하는 물리적 비현실성이 발생할 수 있어 \chg{첫 번째 요구조건인 경계
보장을 충족하지 못한다}.

\textbf{클리핑(hard clipping)}은 $f(z) = \min(\max(z, z_{\min}), z_{\max})$로서
입력값이 경계 $[z_{\min}, z_{\max}]$ 내이면 그대로 통과시키고 경계를 초과하면 강제로
절단하는 경성 바운딩 함수이다. 클리핑은 경계 조건을 강제로 보장하나, 경계 지점에서
도함수가 불연속이 되고 \chg{경계 외의 다수 입력값이 단일 경계값으로 매핑되면서} 경계
근방의 미세한 입력 변화가 출력에 전혀 반영되지 않는 정보 손실이 발생한다.
\chg{즉 클리핑은 경계 보장 요구조건은 충족하지만, 경계 지점에서 매끄럽지 않아 평균 영역의
선형성과 극단 영역의 점진적 포화를 함께 만족하지 못하며, 경계 외에서 입력 변화가 출력에
반영되지 않아 단조성도 부분적으로 훼손된다.}

\textbf{시그모이드 바운딩}은 매끄러운 S자 곡선을 통해 출력을 유한 범위 내로
자연스럽게 제한하며, 전 구간에서 미분 가능하여 입력 변수의 변화에 대한
출력의 민감도를 연속적으로 분석할 수 있다.
\chg{시그모이드 바운딩은 매끄러운 유한 제한으로 경계를 보장하고, 원점 부근에서 근사적으로
선형이며, 입력이 커질수록 출력이 점진적으로 포화하고, 두 채널을 각각의 물리적 범위에서
독립적으로 바운딩할 수 있으며, 전 구간에서 단조증가한다. 따라서 앞서 제시한 다섯 가지
설계 요구조건을 모두 충족하는 것은 세 후보 중 시그모이드 바운딩뿐이며, 본 연구가 수정
시그모이드를 채택한 근거가 여기에 있다.}

\p{[}표~\ref{tab:bounding-comparison}\p{]}에 세 기법의 특성을 종합적으로 비교하였다.

\begin{table}[htbp]
\centering
\caption{바운딩 기법의 비교}
\label{tab:bounding-comparison}
\small
\begin{tabularx}{\textwidth}{@{}>{\raggedright\arraybackslash}p{2.4cm}
  >{\raggedright\arraybackslash}X
  >{\raggedright\arraybackslash}X
  >{\raggedright\arraybackslash}X@{}}
\toprule
\textbf{특성} & \textbf{선형 모델} & \textbf{클리핑(경성)} & \textbf{시그모이드 바운딩} \\
\midrule
수식 &
$R_0 + \alpha x$ &
$\min(\max(\cdot, 0), R_{\max})$ &
$R_0 + E_{O,\text{bnd}} + E_{U,\text{bnd}}$ \\
\addlinespace
경계 보장 &
미보장 (음수/초과 가능) &
보장 (강제 절단) &
보장 (점근적 접근) \\
\addlinespace
미분 가능성 &
전 구간 ($C^{\infty}$) &
경계점 불연속 &
전 구간 ($C^{\infty}$) \\
\addlinespace
경계 정보 보존 &
해당 없음 &
경계에서 기울기 $= 0$ (정보 손실) &
경계에서도 기울기 $> 0$ (정보 보존) \\
\addlinespace
포화 효과 &
없음 (무한 선형 증가) &
경성 포화 (즉시 절단) &
연성 포화 (점진적 둔화) \\
\addlinespace
물리적 해석 &
비현실적 극단값 가능 &
현실적이나 불연속 &
현실적이고 연속적 \\
\bottomrule
\end{tabularx}
\end{table}

\figref{fig:bounding-comparison}\refpo{fig:bounding-comparison}{은}{는} 세 바운딩 기법의 출력 특성을 비교한 것이다.
\chg{선형 모델은 경계를 초과할 수 있고 클리핑은 경계점에서 불연속인 반면, 시그모이드는 $(-1, +1)$ 범위 안에서 경계에 매끄럽게 점근하면서도 전 구간에서 기울기가 양으로 유지되는 점근적 바운딩을 제공한다.}
시그모이드 바운딩은 경계 영역에서의 정보 보존성, 미분 가능성,
운영 해석 가능성 측면에서 $\DRTO$ 프레임워크에 가장 적합하다.

\begin{figure}[htbp]
  \centering
  \begin{tikzpicture}
    \begin{axis}[
      width=0.88\textwidth,
      height=0.52\textwidth,
      xlabel={입력 $z$},
      ylabel={출력},
      xmin=-4.5, xmax=4.5,
      ymin=-1.6, ymax=1.9,
      grid=major,
      grid style={black},
      legend pos=south east,
      legend style={font=\small, rounded corners=2pt, fill=white, fill opacity=0.9},
      axis lines=middle,
      every axis x label/.style={at={(axis description cs:1.02,0.5)}, anchor=west, font=\small},
      every axis y label/.style={at={(axis description cs:0.5,1.02)}, anchor=south, font=\small},
      clip=false
    ]
      % 선형 모델 (제한 없음)
      \addplot[thick, blue!70, domain=-4.5:4.5, samples=100]
        {0.4*x};
      \addlegendentry{선형 ($0.4z$)}

      % 클리핑
      \addplot[thick, red!70, domain=-4.5:4.5, samples=200]
        {min(max(0.4*x, -1), 1)};
      \addlegendentry{클리핑 (hard)}

      % 시그모이드 S(z) = 2/(1+exp(-z)) - 1
      \addplot[thick, black, domain=-4.5:4.5, samples=200]
        {2/(1+exp(-x)) - 1};
      \addlegendentry{시그모이드 $S(z)$}

      % 경계 초과 영역 채우기 (위쪽: x>2.5에서 0.4x > 1)
      \fill[blue!12, opacity=0.5]
        (axis cs:2.5,1) -- (axis cs:4.5,1.8) -- (axis cs:4.5,1) -- cycle;

      % 경계 초과 영역 채우기 (아래쪽: x<-2.5에서 0.4x < -1)
      \fill[blue!12, opacity=0.5]
        (axis cs:-2.5,-1) -- (axis cs:-4.5,-1.8) -- (axis cs:-4.5,-1) -- cycle;

      % 경계선 표시
      \addplot[thin, dashed, black] coordinates {(-4.5, 1) (4.5, 1)};
      \addplot[thin, dashed, black] coordinates {(-4.5, -1) (4.5, -1)};

      % 경계 라벨
      \node[font=\scriptsize, text=black, anchor=west] at (axis cs:4.6, 1) {$+1$};
      \node[font=\scriptsize, text=black, anchor=west] at (axis cs:4.6, -1) {$-1$};

      % 선형 경계 초과 영역 라벨 (위쪽)
      \node[font=\scriptsize, text=black, fill=white, inner sep=1.5pt,
        anchor=south west] at (axis cs:2.8, 1.45)
        {선형: 경계 초과};
      \draw[-{Stealth[length=1.5mm]}, blue!70, thin]
        (axis cs:3.2, 1.42) -- (axis cs:3.5, 1.15);

      % 선형 경계 초과 영역 라벨 (아래쪽)
      \node[font=\scriptsize, text=black, fill=white, inner sep=1.5pt,
        anchor=north east] at (axis cs:-2.8, -1.45)
        {선형: 경계 초과};
      \draw[-{Stealth[length=1.5mm]}, blue!70, thin]
        (axis cs:-3.2, -1.42) -- (axis cs:-3.5, -1.15);

    \end{axis}
  \end{tikzpicture}
  \caption{바운딩 기법 비교}
  \label{fig:bounding-comparison}
\end{figure}


%% -------------------------------------------------------------------
\subsection{개별 시그모이드 바운딩의 적용}
\label{subsec:individual-sigmoid}
%% -------------------------------------------------------------------

본 연구에서 채택한 개별 시그모이드 바운딩(individual sigmoid bounding) 모델은
관측성과 불확실성의 효과를 각각 고유한 물리적 범위(span)에서
독립적으로 바운딩하는 구조이다. 대안인 결합(joint) 시그모이드 $S(z_O + z_U)$가
두 효과를 단일 함수로 압축하는 반면, 개별 시그모이드는 각 효과의 물리적 범위를
독립적으로 보장한다.
%% [G19/일원화] 구 핵심수식(eq:drto-individual-sigmoid, eq:lit-E_O_bnd, eq:lit-E_U_bnd, eq:lit-z_O, eq:lit-z_U)은
%% 2.7절 모델 구조(eq:drto, eq:E_O_bnd, eq:E_U_bnd, eq:z_O, eq:z_U; canonical)와 완전 중복 → 수식 삭제·전방참조(라벨 참조 0건 확인).
\chg{구체적으로 최종 $\DRTO$는 기준선 $R_0$에 관측성 바운딩 조정량
$E_{O,\text{bnd}} = -R_0 \cdot S(z_O)$(물리적 범위 $R_0$)와 불확실성 바운딩 조정량
$E_{U,\text{bnd}} = (R_{\max} - R_0) \cdot S(z_U)$(물리적 범위 $R_{\max} - R_0$)를 가산하여
산출한다. 음수 부호 $-R_0$는 관측성이 높을수록 복구 시간 단축 효과로 작용함을, $z_U$의
분모 $R_{\max} - R_0$는 불확실성 효과를 가용 상향 여유 범위로 정규화함을 뜻한다. 시그모이드
입력 인자 $z_O$, $z_U$의 전체 정식과 변수 및 매개변수의 형식적 정의, 경계 보장의 엄밀한
증명, 곡률 $\kappa = 1.2$의 선정 근거는 모델 구조와 최적화를 다루는 2.7절과 2.10절에서
종합하므로, 본 절은 그 형태만 바운딩 기법의 귀결로 제시한다.}

\chg{이 구조는 세 가지 핵심적 장점을 갖는다. 첫째, 경계 보장의 보편성으로서, 관측성 효과의
바운딩 범위($R_0$)와 불확실성 효과의 바운딩 범위($R_{\max} - R_0$)를 각각의 물리적 의미에
맞게 독립적으로 설정하므로 $\DRTO \in [0,\; R_{\max}]$가 $\kappa$ 값에 무관하게 수학적으로
보장된다. 둘째, 경사 최적화와 민감도 분석의 용이성으로서, 각 시그모이드 함수가 전체
정의역에서 $C^{\infty}$이므로 경사 기반 최적화와 민감도 분석이 용이하다. 셋째, C1 조건
분리 가능성으로서, C1 조건($E_{U,\text{bnd}} = 0$)에서 관측성만의 순수 효과를 불확실성
상한과 독립적으로 분석할 수 있어 모델의 단계적 확장이 자연스럽다.}


%% -------------------------------------------------------------------
\subsection{이중채널 감쇠 메커니즘}
\label{subsec:dual-channel-attenuation}
%% -------------------------------------------------------------------

개별 시그모이드 바운딩 구조는 관측성 채널($E_{O,\text{bnd}} = -R_0 \cdot S(z_O)$,
복구 시간 단축)과 불확실성 채널($E_{U,\text{bnd}} = (R_{\max}-R_0) \cdot S(z_U)$,
복구 시간 연장)이 동시에 작용하는 이중채널 구조를 형성한다. 두 채널이 각자의 물리적
범위에서 독립적으로 바운딩되므로, 불확실성이 극단적으로 커지는 꼬리(tail) 영역에서는
불확실성 채널이 지배적으로 작용하면서 관측성 채널의 단축 효과가 구조적으로 약화되는
현상이 나타난다. 이를 이중채널 감쇠(dual-channel attenuation)라 한다.

이 메커니즘은 실무적으로 중요한 함의를 갖는다. 고도의 모니터링 체계를 갖춘
조직(높은 관측성)이라 하더라도 대규모 사이버 공격과 같은 극단적 불확실성 상황에서는
관측성 투자만으로 복구 시간을 충분히 단축하기 어려우며, 불확실성 자체에 대한 관리,
이를테면 업무연속성계획(BCP) 시나리오의 다양화나 사이버 보험 등이 병행되어야 함을
시사한다. 이중채널 감쇠가 분포 수준에서 발현되는 전환점과 그 정량적 크기(꼬리
영역에서의 관측성 효과 감쇠 비율, 조건부 위험가치 격차 등)는 본 연구의 실증 결과에서
분할 회귀를 통해 확인되며, 이는 \secref{sec:ch5_dual_channel}에서 상세히 분석한다.
